% problem_id: S001-theorem-curves-rational-maps
% source_id: S001 (Stacks Project)
% source_file: curves.tex
% statement_locator: curves.tex lines 158--170
% proof_locator: curves.tex lines 172--228
% env: theorem
% id_note: label 在本来源内唯一
% status: complete (statement + author proof verbatim + reason + locators)
%
\begin{theorem}
\label{theorem-curves-rational-maps}
Let $k$ be a field. The following categories are canonically equivalent
\begin{enumerate}
\item The category of finitely generated field extensions $K/k$ of
transcendence degree $1$.
\item The category of curves and dominant rational maps.
\item The category of normal projective curves and nonconstant morphisms.
\item The category of nonsingular projective curves and nonconstant morphisms.
\item The category of regular projective curves and nonconstant morphisms.
\item The category of normal proper curves and nonconstant morphisms.
\end{enumerate}
\end{theorem}

\begin{proof}
The equivalence between categories (1) and (2) is the restriction of the
equivalence of
Varieties, Theorem \ref{varieties-theorem-varieties-rational-maps}.
Namely, a variety is a curve if and only if its function field has
transcendence degree $1$, see for example
Varieties, Lemma \ref{varieties-lemma-dimension-locally-algebraic}.

\medskip\noindent
The categories in (3), (4), (5), and (6) are the same. First of all, the
terms ``regular'' and ``nonsingular'' are synonyms, see
Properties, Definition \ref{properties-definition-regular}.
Being normal and regular are the same thing for Noetherian
$1$-dimensional schemes
(Properties, Lemmas \ref{properties-lemma-regular-normal} and
\ref{properties-lemma-normal-dimension-1-regular}). See
Varieties, Lemma \ref{varieties-lemma-regular-point-on-curve}
for the case of curves. Thus (3) is the same as (5). Finally, (6)
is the same as (3) by
Varieties, Lemma \ref{varieties-lemma-dim-1-proper-projective}.

\medskip\noindent
If $f : X \to Y$ is a nonconstant morphism of nonsingular projective curves,
then $f$ sends the generic point $\eta$ of $X$ to the generic point $\xi$ of
$Y$. Hence we obtain a morphism
$k(Y) = \mathcal{O}_{Y, \xi} \to \mathcal{O}_{X, \eta} = k(X)$
in the category (1). If two morphisms $f,g: X \to Y$ gives the same morphism
$k(Y) \to k(X)$, then by the equivalence between (1) and (2),
$f$ and $g$ are equivalent as rational maps, so $f=g$ by
Lemma \ref{lemma-extend-over-normal-curve}.
Conversely, suppose that we have a map
$k(Y) \to k(X)$ in the category (1). Then we obtain a morphism $U \to Y$
for some nonempty open $U \subset X$. By Lemma \ref{lemma-extend-over-dvr}
this extends to all of $X$ and we obtain a morphism in the category (5).
Thus we see that there is a fully faithful functor (5)$\to$(1).

\medskip\noindent
To finish the proof we have to show that every $K/k$ in (1)
is the function field of a normal projective curve.
We already know that $K = k(X)$ for some curve $X$.
After replacing $X$ by its normalization
(which is a variety birational to $X$)
we may assume $X$ is normal
(Varieties, Lemma \ref{varieties-lemma-normalization-locally-algebraic}).
Then we choose $X \to \overline{X}$ with
$\overline{X} \setminus X = \{x_1, \ldots, x_n\}$ as in
Varieties, Lemma \ref{varieties-lemma-reduced-dim-1-projective-completion}.
Since $X$ is normal and since each
of the local rings $\mathcal{O}_{\overline{X}, x_i}$ is normal
we conclude that $\overline{X}$ is a normal projective curve as desired.
(Remark: We can also first compactify using
Varieties, Lemma \ref{varieties-lemma-dim-1-projective-completion}
and then normalize using
Varieties, Lemma \ref{varieties-lemma-normalization-locally-algebraic}.
Doing it this way we avoid using the somewhat tricky
Morphisms, Lemma \ref{morphisms-lemma-relative-normalization-normal-codim-1}.)
\end{proof}
